\documentclass[twocolumn]{article}
\input{preambulo.tex}

\begin{document}
\thispagestyle{firstpage}

\twocolumn[
\articlefront
  {Working draft, 27 October 2026}
  {Detecção de lesões periapicais em radiografias panorâmicas com poucos exemplos locais: Random Forest por pixel versus U-Net treinada em dados públicos}
  {William Jota\textsuperscript{1}}
  {Lesões periapicais aparecem em radiografias panorâmicas como áreas radiolúcidas pequenas e de contraste baixo, e parte delas não é registrada na leitura de rotina. Avaliamos se um detector automático poderia ser construído para o acervo de uma clínica que dispunha de apenas 20 panorâmicas com lesões contornadas por especialista (45 lesões) e de 422 panorâmicas sem marcação. Comparamos um classificador Random Forest por pixel, com 38 atributos de intensidade, textura e contexto, e uma U-Net com codificador ResNet34. Ambos foram treinados com um conjunto público de 3.924 panorâmicas com lesões delimitadas por caixas, convertidas em elipses. Cada paciente local foi avaliado por um modelo que não o viu no treinamento, e uma lesão foi considerada encontrada quando uma região prevista tinha ao menos 20\% de sua área sobre ela. O Random Forest encontrou no máximo 16\% das lesões (7/45; IC95\% 8--29\%) e gerou 5,8 falsos positivos por imagem nesse ponto, porque respondeu à textura do osso trabecular e à posição na imagem. A U-Net encontrou 64\% das lesões (29/45; IC95\% 50--77\%) sem falsos positivos nas 20 imagens (precisão 29/29; IC95\% 88--100\%) e Dice de 0,57. O ajuste fino com as imagens locais acrescentou uma lesão (67\%) e elevou o Dice a 0,60. As lesões não encontradas eram menores (diâmetro equivalente mediano de 20 contra 41 px) e concentradas na maxila (50\% encontradas, contra 76\% na mandíbula). Aplicada às 422 panorâmicas, a U-Net assinalou 673 lesões em 376 imagens. Esses resultados indicam que dados públicos rotulados de forma grosseira bastam para uma rede convolucional detectar lesões periapicais evidentes em outro serviço, mas não para um classificador por pixel. Lesões pequenas e maxilares, a precisão no acervo sem marcação e a concordância entre examinadores ainda precisam de validação prospectiva.}
  {\textbf{Palavras-chave:} lesão periapical; periodontite apical; radiografia panorâmica; segmentação; U-Net; Random Forest; aprendizado por transferência.}
  {\textsuperscript{1} Laboratório de Estudos Eletromagnéticos (LEMAG), Universidade de Brasília, Brasília, Brasil.}
]

\section{Introdução}

A periodontite apical é a resposta inflamatória dos tecidos ao redor do ápice radicular à infecção do sistema de canais, e sua manifestação radiográfica típica é uma área radiolúcida junto à ponta da raiz. Em escala populacional, a radiografia panorâmica é o exame mais disponível para esse rastreamento, mas sobreposição de estruturas, distorção e baixo contraste reduzem sua sensibilidade em relação à tomografia computadorizada de feixe cônico [1]. Mesmo especialistas deixam passar parte das lesões: em uma comparação com 24 cirurgiões bucomaxilofaciais, a taxa média de verdadeiros positivos na leitura de panorâmicas foi de 0,51 [2].

Redes neurais convolucionais vêm sendo usadas para apoiar essa leitura. Ekert e colaboradores mostraram detecção de lesões apicais em panorâmicas com uma rede convolucional [3], e Endres e colaboradores relataram desempenho comparável ao de parte dos especialistas com 2.902 panorâmicas rotuladas [2]. Esses estudos, porém, dependem de milhares de imagens marcadas no próprio serviço. Clínicas e grupos de pesquisa menores raramente dispõem desse volume, e a marcação densa de lesões por especialistas é cara.

Duas alternativas se apresentam a quem tem poucos exemplos locais. A primeira é a classificação por pixel com atributos predefinidos e Random Forest [4], popularizada por ferramentas interativas como ilastik e Trainable Weka Segmentation [5,6], que aprendem com poucos exemplos e são interpretáveis. A segunda é treinar uma rede convolucional em um conjunto público de outro serviço e aplicá-la ao acervo local. Conjuntos públicos de panorâmicas com lesões periapicais já existem [7,8], mas suas marcações costumam ser caixas ou contornos do dente inteiro, e não o contorno da lesão.

O problema prático deste estudo é decidir qual dessas vias funciona quando o acervo local tem 20 panorâmicas contornadas e 422 sem marcação. Comparamos um Random Forest por pixel e uma U-Net [9], ambos treinados em 3.924 panorâmicas públicas de Hanói [7], sob o mesmo protocolo de avaliação por lesão nas imagens locais. Também medimos o efeito de um ajuste fino com as imagens locais, caracterizamos que lesões cada modelo deixa de encontrar e aplicamos o modelo escolhido ao acervo sem marcação. A contribuição é uma comparação justa entre as duas vias no cenário de poucos dados, acompanhada de uma delimitação explícita do que a validação disponível permite afirmar.

\widefigure{95mm}{fig3_diagrama.pdf}{\textbf{Fluxo de trabalho.} Dados, pré-processamento, modelos, validação e aplicação. Somente o conjunto público de Hanói foi enviado ao ambiente de nuvem para o treinamento da U-Net; as radiografias de pacientes foram processadas apenas no computador local.}

\section{Materiais e métodos}

\subsection{Radiografias locais e marcações}

O acervo local contém 440 radiografias panorâmicas digitais em tons de cinza (438 em JPEG e 2 em TIFF), com largura entre cerca de 2.400 e 2.900 pixels. Vinte pacientes tiveram as lesões periapicais contornadas por \pendente{número e formação dos examinadores} em um documento PDF, uma página por paciente; dezoito dessas radiografias também estão no acervo e duas existem apenas no PDF. As 422 radiografias restantes, sem marcação, formaram o conjunto de aplicação.

Os contornos estavam gravados no PDF como desenhos vetoriais sobre a radiografia. Reconstruímos cada radiografia unindo as faixas de imagem da página e rasterizamos os contornos nas mesmas coordenadas, o que dispensou registro entre imagem e marcação. Muitos traços eram arcos abertos em forma de ``$\cap$'', contornando a lesão e excluindo a raiz. Fechamos cada arco ligando suas extremidades livres e preenchemos o interior; quando não havia interior, usamos o fecho convexo do traço. Componentes com menos de 12 pixels de lado foram descartados como resquícios do traço. O resultado foram 45 lesões em 20 pacientes (Fig.~2a), 20 delas na metade superior da imagem, tratada aqui como maxila, e 25 na inferior, tratada como mandíbula.

\subsection{Conjunto público de treinamento}

Usamos o conjunto de panorâmicas com lesões periapicais da Universidade de Medicina de Hanói [7], distribuído sob licença CC BY 4.0 e composto por 3.924 radiografias originais. Cada lesão está delimitada por uma caixa justa ao redor da área radiolúcida, com largura mediana de 60 pixels, e recebe uma classe 3, 4 ou 5. Essa classe pode corresponder ao índice periapical [10], mas a documentação não o confirma, e ela não foi usada. Como a caixa contém lesão e osso adjacente, cada caixa foi convertida em uma elipse inscrita com semieixos a 85\% da meia-caixa. Os pixels entre essa elipse e uma elipse a 150\% foram excluídos do treinamento, porque podem ser lesão ou fundo (Fig.~2b).

Também examinamos o conjunto DENTEX [8], que contém 116 panorâmicas com lesão periapical. Nele, porém, o polígono de cada achado contorna o dente inteiro, e não a lesão, e o conjunto não foi usado.

\begin{figure*}[!t]
\centering
\includegraphics[width=\textwidth]{fig2_rotulos.png}
\caption{\textbf{Rótulos usados no treinamento e na avaliação.} \textbf{a}, marcação de um especialista no PDF (traço) e a máscara derivada após o fechamento do arco e o preenchimento; a raiz contornada pelo traço fica fora da máscara. \textbf{b}, panorâmica do conjunto de Hanói com a caixa convertida em elipse (lesão) e o anel ao redor excluído do treinamento.}
\end{figure*}

\subsection{Pré-processamento}

Todas as radiografias foram reduzidas a 1.400 pixels de largura, preservando a proporção, por interpolação de área, e tiveram o contraste equalizado por CLAHE (limite 2,0; grade de $8 \times 8$) [11]. Essa etapa aproxima o contraste de aparelhos diferentes. Nessa escala, a lesão marcada típica tem diâmetro equivalente mediano de 32 pixels (intervalo interquartil 20--45). Medidas de tamanho são dadas em pixels dessa escala, porque as imagens não continham calibração espacial.

\subsection{Random Forest por pixel}

Cada pixel foi descrito por 38 atributos: intensidade; suavizações gaussianas ($\sigma = 1, 2, 4, 8, 16$ px); diferenças de gaussianas (1--4, 2--8, 4--16, 8--32 e 32--64 px); magnitude do gradiente ($\sigma = 1, 2, 4, 8$); autovalores da Hessiana ($\sigma = 2, 4, 8$); desvio-padrão local ($\sigma = 4, 8$); contexto, dado pela diferença entre a intensidade suavizada ($\sigma = 4, 8$) a 20 e 40 px acima e abaixo e a 25 px para cada lado e a do próprio pixel; e posição horizontal, vertical e distância horizontal ao centro. Os atributos de contexto foram incluídos porque a lesão é definida pela relação com a raiz adjacente.

Sorteamos 200 pixels de lesão e 1.000 de fundo em cada uma de 800 panorâmicas de Hanói (959.989 pixels; 16,7\% de lesão). Metade do fundo veio da faixa central da imagem, onde estão raízes e espaços interdentais. O Random Forest tinha 100 árvores, profundidade máxima 20, mínimo de 10 amostras por folha, $\sqrt{d}$ atributos por divisão e 200.000 pixels por árvore. Esse modelo foi chamado RF B1. Em seguida, aplicamos o RF B1 a outras 300 panorâmicas de Hanói e acrescentamos até 600 pixels de fundo por imagem onde a probabilidade de lesão era $\geq 0{,}3$ (180.000 negativos difíceis). O modelo retreinado com esses exemplos foi chamado RF B. A variante RF C somou ao RF B, em validação cruzada, 2.000 pixels de lesão e 8.000 de fundo de cada imagem local de treino.

\subsection{U-Net}

A U-Net [9] usou como codificador uma ResNet34 [12] pré-treinada no ImageNet, com a imagem em tons de cinza replicada nos três canais, e um decodificador com cinco estágios de duas convoluções $3 \times 3$, normalização em lote e conexões de atalho. O treinamento usou 3.724 panorâmicas de Hanói, reservando 200 para validação interna. A cada passo, um lote de 8 recortes de $512 \times 768$ pixels foi sorteado; metade foi centrada em uma lesão, para que ela não fosse rara demais. Os recortes receberam espelhamento horizontal e variações de gama, ganho e brilho. A perda foi a soma da entropia cruzada binária com a perda Dice, ambas calculadas somente fora do anel excluído. Usamos AdamW [13] com taxa de aprendizado $3 \times 10^{-4}$, decaimento $10^{-4}$, aquecimento de 200 passos e decaimento cossenoidal ao longo de 6.000 passos, em precisão mista, em uma GPU NVIDIA T4 do Google Colab (cerca de 0,45 s por passo). Salvamos o modelo de maior Dice na validação interna. Na previsão, a probabilidade foi a média entre a imagem e sua versão espelhada.

\subsection{Ajuste fino com as imagens locais}

Para testar se as 20 imagens locais acrescentavam informação, ajustamos a U-Net com os contornos dos especialistas. Para caber em CPU e limitar a memorização de poucos exemplos, congelamos o codificador e as estatísticas de normalização e ajustamos apenas o decodificador. Usamos 150 passos com lotes de 2 recortes de $384 \times 576$ pixels, 70\% deles centrados em lesão, e AdamW com taxa $3 \times 10^{-4}$. Uma faixa de 4 pixels em cada lado do contorno marcado foi excluída da perda, porque o traço manual não é exato.

\subsection{Detecção e métricas por lesão}

O mapa de probabilidade de cada modelo foi convertido em lesões pelo mesmo procedimento. O mapa foi suavizado (gaussiana, $\sigma = 4$ px), e cada máximo local em uma janela de $31 \times 31$ pixels com probabilidade $\geq 0{,}05$ originou uma lesão candidata. A região da candidata foi o conjunto conectado de pixels acima de metade do pico, limitado a 40 pixels do pico para o Random Forest e a 100 pixels para a U-Net, cujo mapa é mais nítido. A confiança de cada candidata foi o valor do pico.

Uma região prevista foi considerada correta quando ao menos 20\% de sua área caía sobre uma lesão marcada; uma lesão marcada foi considerada encontrada quando tocada por uma região correta. Esse critério impede que regiões extensas sejam contadas como acerto apenas por encostarem em uma lesão. Variando o limiar de confiança de 0,05 a 0,95, traçamos curvas FROC, isto é, a sensibilidade por lesão em função dos falsos positivos por imagem [14]. Para cada modelo, o limiar de operação foi o de maior F1 entre sensibilidade e precisão. Como esse limiar foi escolhido nas mesmas 20 imagens, os valores no ponto de operação são levemente otimistas. Relatamos também o Dice entre os pixels previstos e marcados, somados sobre as imagens. Intervalos de confiança de 95\% foram calculados pelo método de Wilson para proporções [15] e por \textit{bootstrap} por paciente (2.000 reamostragens) para o Dice.

\subsection{Protocolo de validação}

O RF B1, o RF B e a U-Net não usaram nenhuma das 20 imagens locais no treinamento; para eles, as 20 imagens são inteiramente de teste. O RF C e a U-Net com ajuste fino usaram as imagens locais e foram avaliados por validação cruzada em cinco partes fixas de quatro pacientes: cada paciente foi previsto por um modelo treinado sem ele.

\subsection{Aplicação e ferramenta de revisão}

O modelo escolhido foi aplicado às 422 panorâmicas sem marcação. Para cada imagem, gravamos a máscara na resolução original, a radiografia com as lesões contornadas e uma tabela com posição, tamanho e confiança de cada lesão. Para medir a precisão nesse conjunto, desenvolvemos uma ferramenta local de revisão. Nela, o especialista confirma, rejeita ou marca como duvidosa cada detecção, desenha as lesões não detectadas e registra imagens inadequadas. A ferramenta permite ocultar as detecções para que a leitura não seja induzida pelo modelo. As revisões geram estimativas de precisão e sensibilidade com IC95\% e novos rótulos de treinamento; nenhuma revisão havia sido concluída na redação deste rascunho.

\subsection{Software e privacidade}

O pré-processamento, o Random Forest e a avaliação foram implementados em Python 3.11 com NumPy, SciPy, OpenCV, scikit-image e scikit-learn 1.5 [16]; a U-Net, em PyTorch [17] e torchvision. As figuras quantitativas foram produzidas em R 4.6 com ggplot2 e patchwork. Somente o conjunto público de Hanói foi enviado ao Google Colab. As radiografias de pacientes foram processadas, avaliadas e revisadas exclusivamente no computador local, em CPU de quatro núcleos (cerca de 7,6 s por panorâmica com a U-Net e 4,7 s com o Random Forest).

\section{Resultados}

\subsection{A U-Net encontrou quatro vezes mais lesões que o Random Forest, sem falsos positivos}

Para decidir qual das duas vias funciona no acervo local, comparamos as curvas FROC dos modelos nas 45 lesões (Fig.~3a, Tabela~1). O melhor Random Forest encontrou no máximo 16\% das lesões (RF B; 7/45; IC95\% 8--29\%), com 5,8 falsos positivos por imagem nesse ponto. Restringindo-se a até 2 falsos positivos por imagem, nenhuma variante passou de 11\%. A inclusão dos negativos difíceis não mudou esse quadro (RF B1: 7\% com 2 falsos positivos por imagem; RF B: 7\%). Também não o mudou a inclusão das imagens locais no treinamento (RF C: 11\% de sensibilidade, 25\% de precisão e 0,6 falso positivo por imagem).

A U-Net treinada apenas com Hanói encontrou 64\% das lesões no ponto de operação (29/45; IC95\% 50--77\%) e 69\% (31/45) no limiar mais permissivo, sem gerar falsos positivos em nenhuma das 20 imagens (precisão 29/29; IC95\% 88--100\%). O Dice foi 0,57 (IC95\% 0,52--0,62), contra 0,14 do RF C. A curva da U-Net termina próxima de zero falso positivo porque, abaixo do limiar de operação, quase não surgem novas candidatas: as lesões não encontradas recebem probabilidade próxima de zero.

\widefigure{140mm}{fig1_validacao.pdf}{\textbf{Validação e escolha do modelo nas 20 panorâmicas locais (45 lesões).} \textbf{a}, curvas FROC; faixas, IC95\% de Wilson ($n = 45$); pontos, limiar de operação (melhor F1). \textbf{b}, sensibilidade, precisão e Dice no limiar de operação; barras de erro, IC95\% (Wilson para proporções; \textit{bootstrap} por paciente para o Dice). Cores como em \textbf{a}. \textbf{c}, lesões marcadas segundo o diâmetro equivalente (px, imagem com 1.400 px de largura) e a metade da imagem, encontradas (preenchidas) ou não (abertas) pela U-Net de Hanói; tracejados, medianas; dispersão vertical apenas para visualização. \textbf{d}, treinamento da U-Net: perda média a cada 50 passos e Dice na validação interna de Hanói (200 imagens) a cada 500 passos; linha vertical, passo do modelo salvo (5.500).}

\begin{table*}[!t]
\caption{\textbf{Desempenho por lesão nas 20 panorâmicas locais (45 lesões).} Sensibilidade com no máximo 0,5, 1, 2 e 4 falsos positivos (FP) por imagem e, à direita, valores no limiar de operação (melhor F1), com IC95\% de Wilson ou \textit{bootstrap}. RF C e U-Net com ajuste fino foram avaliados por validação cruzada em cinco partes; os demais não usaram as imagens locais.}
\centering\tablefont
\begin{tabularx}{\textwidth}{L *{4}{c} Y Y c Y}
\toprule
Modelo & 0,5 FP & 1 FP & 2 FP & 4 FP & Sensibilidade & Precisão & FP/imagem & Dice \\
\midrule
RF B (Hanói) & 4\% & 4\% & 7\% & 11\% & 16\% (8--29) & 6\% (3--11) & 5,8 & 0,04 (0,02--0,06) \\
RF C (Hanói + locais) & 4\% & 11\% & 11\% & 11\% & 11\% (5--23) & 25\% (10--49) & 0,6 & 0,14 (0,06--0,26) \\
U-Net (Hanói) & \textbf{69\%} & \textbf{69\%} & \textbf{69\%} & \textbf{69\%} & 64\% (50--77) & \textbf{100\% (88--100)} & \textbf{0,0} & 0,57 (0,52--0,62) \\
U-Net + ajuste fino & \textbf{69\%} & \textbf{69\%} & \textbf{69\%} & \textbf{69\%} & \textbf{67\% (52--79)} & \textbf{100\% (89--100)} & \textbf{0,0} & \textbf{0,60 (0,55--0,67)} \\
\bottomrule
\end{tabularx}
\end{table*}

\subsection{O Random Forest respondeu ao osso trabecular e à posição}

Os mapas de probabilidade explicam a diferença (Fig.~4). O Random Forest atribuiu probabilidade intermediária a grandes extensões de osso trabecular, sobretudo no corpo da mandíbula, e as regiões resultantes raramente coincidiam com as lesões. Os três atributos de maior importância no RF C foram a posição vertical, a distância horizontal ao centro e a posição horizontal, que somaram 31\% da importância por redução de impureza de Gini. Esse padrão indica que o classificador aprendeu onde as lesões costumam estar, e não como elas se apresentam. A U-Net, ao contrário, produziu mapas concentrados em poucos focos junto aos ápices radiculares.

\begin{figure*}[!t]
\centering
\includegraphics[width=\textwidth]{fig4_rf_vs_unet.png}
\caption{\textbf{Mapas de probabilidade e detecções do Random Forest e da U-Net em três pacientes.} Primeira coluna: radiografia com as lesões marcadas pelo especialista (verde). Segunda e terceira colunas: probabilidade de lesão (azul, de transparente em 0 a opaco em 1), lesões marcadas (verde) e detecções no limiar de operação (ciano). O RF C foi previsto por modelos que não viram o paciente; a U-Net não viu nenhuma imagem local. Em \textbf{b}, as quatro lesões da metade superior da imagem não foram encontradas por nenhum dos modelos.}
\end{figure*}

\subsection{O ajuste fino com as imagens locais teve efeito pequeno}

O ajuste fino da U-Net com as 20 imagens locais encontrou uma lesão a mais no ponto de operação (30/45; 67\%) e elevou o Dice de 0,57 para 0,60, mantendo a ausência de falsos positivos (Tabela~1). A sensibilidade máxima permaneceu em 69\% (31/45), e os intervalos de confiança das duas versões se sobrepõem amplamente. Com 45 lesões, uma lesão de diferença não distingue os modelos. Por isso, a aplicação usou a U-Net sem ajuste fino, que não depende das imagens locais e cuja avaliação é inteiramente independente delas.

\subsection{As lesões não encontradas eram pequenas e maxilares}

Para entender o limite de 69\%, examinamos as 16 lesões que a U-Net não encontrou no ponto de operação (Fig.~3c, Fig.~5). A U-Net encontrou 76\% das lesões da mandíbula (19/25), mas apenas 50\% das da maxila (10/20). As lesões não encontradas tinham diâmetro equivalente mediano de 20 pixels, contra 41 pixels das encontradas. Em 10 dos 20 pacientes, todas as lesões marcadas foram encontradas.

\begin{figure*}[p]
\centering
\includegraphics[width=\textwidth,height=215mm,keepaspectratio]{fig5_validacao_20_pacientes.png}
\caption{\textbf{Detecções da U-Net nas 20 panorâmicas locais.} Verde: lesões marcadas pelo especialista; ciano: detecções da U-Net de Hanói no limiar de operação. O título de cada painel indica o paciente, as lesões encontradas e os falsos positivos, quando houver. Nenhuma das imagens foi usada no treinamento da U-Net.}
\end{figure*}

\begin{table*}[!t]
\caption{\textbf{Análises de sensibilidade da U-Net nas 20 panorâmicas locais.} Os limiares e as combinações foram examinados nas mesmas imagens; os resultados descrevem o comportamento do modelo e não foram usados para escolher a versão aplicada.}
\centering\tablefont
\begin{tabularx}{\textwidth}{>{\RaggedRight\arraybackslash}p{52mm} c c L}
\toprule
Variante & Lesões encontradas & FP/imagem & Leitura \\
\midrule
Limiar de operação (0,5) & 29/45 (64\%) & 0,00 & referência \\
Limiar 0,2 & 31/45 (69\%) & 0,20 & duas lesões a mais \\
Limiar e pico mínimo 0,002 & 33/45 (73\%) & 0,75 & ganho pequeno, com quase um falso positivo por imagem \\
Imagem e ampliação 1,5$\times$ (máximo) & 33/45 (73\%) & $\leq$1,0 & F1 0,76, contra 0,78 na escala original \\
Somente ampliação 1,5$\times$ & 28/45 (62\%) & $\leq$1,0 & piora \\
Somente ampliação 2$\times$ & 3/45 (7\%) & $\leq$0,5 & a rede não reconhece lesões fora da escala de treino \\
\bottomrule
\end{tabularx}
\end{table*}

Reduzir o limiar não recuperou essas lesões. Mesmo aceitando qualquer pico acima de 0,002, a sensibilidade chegou a 73\%, com 0,75 falso positivo por imagem (Tabela~2). Prever também na imagem ampliada 1,5 vez, para tornar as lesões pequenas maiores, encontrou duas lesões adicionais somente ao custo de cerca de um falso positivo por imagem. A ampliação de 2 vezes levou a sensibilidade a 7\%. Esses resultados indicam que as lesões perdidas não estão logo abaixo do limiar: a rede, treinada com lesões de outro tamanho típico, não as reconhece.

\subsection{Aplicação a 422 panorâmicas sem marcação}

A U-Net de Hanói assinalou 673 lesões em 376 das 422 panorâmicas (89\%); 46 imagens não tiveram detecção, e 28 tiveram quatro ou mais (Fig.~6a, Fig.~7). As detecções concentraram-se junto aos ápices das duas arcadas, com 69\% delas na metade inferior da imagem e distribuição equilibrada entre os lados (Fig.~6b). O diâmetro equivalente mediano das lesões detectadas foi 41 pixels (intervalo interquartil 31--53), semelhante ao das lesões encontradas na validação e maior que o das marcadas em geral (32 pixels). Esse padrão sugere que lesões pequenas também foram perdidas nesse conjunto (Fig.~6c). A confiança foi $\geq 0{,}95$ em 86\% das detecções (Fig.~6d) e, por isso, distingue pouco os casos duvidosos. Sem marcação de referência, a precisão nesse conjunto não pode ser medida, e a proporção de 89\% de imagens com detecção ainda precisa ser confirmada pela revisão especializada.

\widefigure{100mm}{fig2_aplicacao.pdf}{\textbf{Aplicação da U-Net às 422 panorâmicas sem marcação.} \textbf{a}, número de lesões detectadas por imagem. \textbf{b}, posição dos centros das 673 detecções (fração da largura e da altura da imagem); manchas, densidade bidimensional por \textit{kernel}; as linhas tracejadas são referências aproximadas, não anatômicas. \textbf{c}, distribuição acumulada do diâmetro equivalente das lesões detectadas na aplicação e das lesões marcadas na validação, encontradas e não encontradas; tracejados, medianas. \textbf{d}, confiança das detecções (pico do mapa de probabilidade suavizado).}

\begin{figure*}[!t]
\centering
\includegraphics[width=\textwidth]{fig7_aplicacao_exemplos.png}
\caption{\textbf{Exemplos de detecções da U-Net no acervo sem marcação.} Uma imagem sorteada em cada faixa de número de detecções (0, 1, 2--3 e 4 ou mais). Os identificadores (A001--A422) são anônimos. As detecções não foram confirmadas por especialista.}
\end{figure*}

\scopenote{\textbf{Escopo da evidência.} As métricas de validação vêm de 45 lesões em 20 panorâmicas de um único serviço, marcadas por \pendente{número de} examinador(es), com limiar de operação escolhido nas mesmas imagens. Não há leitura repetida para medir a variabilidade entre examinadores. Os resultados no acervo de 422 imagens descrevem o que a U-Net assinalou, e não sua precisão. Nenhum resultado deste rascunho é diagnóstico.}

\section{Discussão}

O resultado central é que uma rede convolucional treinada apenas com panorâmicas públicas de outro país, rotuladas por caixas, encontrou dois terços das lesões periapicais marcadas em um serviço local, sem falsos positivos, enquanto um classificador por pixel treinado com os mesmos dados não passou de 16\%. A diferença não se explica pelo volume de dados nem pelo rótulo grosseiro, que eram os mesmos para ambos. Ela decorre do que cada modelo consegue representar: a lesão periapical é uma radiolucidez definida pela relação com uma raiz específica, e não uma textura local. Mesmo com atributos de contexto e negativos difíceis, o Random Forest acabou respondendo à textura trabecular e à posição na imagem, enquanto a U-Net integra o contexto em várias escalas.

Essa conclusão contrasta com o sucesso de classificadores por pixel em outras áreas de imagem, nas quais o alvo tem assinatura local própria, como cor ou textura distintas [5,6]. Aqui, a assinatura local da lesão coincide com a do osso trabecular menos denso, e o classificador por pixel não tem como separá-las sem modelar a anatomia. O Random Forest permanece útil como linha de base transparente, e sua falha documenta por que a escolha da arquitetura importa mais que o ajuste de atributos nesse problema.

A sensibilidade de 64--69\% deve ser lida com dois cuidados. Primeiro, ela foi medida em 45 lesões de 20 pacientes, com intervalo de confiança de 50 a 77\%, e o limiar foi escolhido nas mesmas imagens. Segundo, a referência é a marcação de \pendente{número de} examinador(es) sobre panorâmicas, exame cuja sensibilidade é limitada em relação à tomografia [1]. Sem leitura repetida, não sabemos quantas das lesões não encontradas seriam reconhecidas por outro especialista. A comparação com a taxa média de verdadeiros positivos de 0,51 relatada para cirurgiões em panorâmicas [2] é sugestiva, mas não direta, porque as populações, as referências e os critérios de acerto diferem.

As falhas foram sistemáticas, e não aleatórias. A U-Net perdeu sobretudo lesões pequenas e maxilares, e nem o limiar mais permissivo nem a ampliação da imagem as recuperaram. A interpretação mais simples é que essas lesões estão sub-representadas no conjunto de treinamento ou diferem das lesões de Hanói em tamanho e aparência. Na maxila, a sobreposição do seio maxilar e do palato também reduz o contraste. Essa leitura leva a três intervenções verificáveis: treinar em resolução maior, usar uma perda que penalize mais a lesão não encontrada, como a perda de Tversky [18], e acrescentar exemplos locais de lesões pequenas e maxilares. As duas primeiras estão em teste. A terceira depende da ferramenta de revisão.

O ajuste fino com 20 imagens quase não alterou o desempenho. Esse resultado é compatível com a ideia de que o gargalo não é a adaptação ao aparelho local, já reduzida pela equalização de contraste, mas a falta de exemplos dos casos difíceis. Vinte imagens com 45 lesões, das quais apenas 16 foram perdidas, não bastam para ensinar à rede um padrão que ela não aprendeu em 3.724 imagens. A melhora do Dice, de 0,57 para 0,60, sugere que as imagens locais ajudam sobretudo no contorno, e não na detecção.

Os rótulos de Hanói também limitam o contorno. As elipses inscritas nas caixas fazem a U-Net desenhar regiões arredondadas, que não acompanham lesões irregulares; daí um Dice moderado mesmo quando a lesão é encontrada. Modelos de segmentação guiados por caixa, como o MedSAM [19], poderiam converter as caixas em contornos mais fiéis sem nova marcação manual.

No acervo de 422 imagens, a ausência de referência impede afirmar acurácia. A proporção de 89\% de imagens com ao menos uma detecção pode refletir a população atendida, falsos positivos não observados na validação ou ambos. Como a validação local não teve nenhum falso positivo, a primeira hipótese não pode ser descartada, mas também não pode ser confirmada. A revisão de uma amostra aleatória dessas imagens, e não apenas das prioritárias, é o passo que transforma a aplicação em evidência de desempenho. Duas leituras independentes de um subconjunto comum também medirão a concordância entre examinadores, teto realista para qualquer modelo.

Também há limitações de dados. A separação entre maxila e mandíbula foi aproximada pela metade da imagem. O acervo contém possíveis duplicatas de pacientes com grafias diferentes, e uma das 45 máscaras incluiu a raiz porque o preenchimento automático recorreu ao fecho convexo. Nenhuma dessas limitações altera a comparação entre os modelos, mas todas devem ser resolvidas antes de uma avaliação definitiva.

\section{Conclusão}

Com apenas 20 panorâmicas marcadas no próprio serviço, uma U-Net treinada em 3.924 panorâmicas públicas encontrou 64--69\% das lesões periapicais sem falsos positivos, enquanto um Random Forest por pixel com os mesmos dados não passou de 16\%. O ajuste fino local teve efeito pequeno, e as lesões não encontradas foram sistematicamente pequenas e maxilares. A evidência atual justifica usar a U-Net como apoio à triagem de lesões evidentes, sob revisão especializada. Ela não justifica afirmar desempenho diagnóstico, precisão no acervo sem marcação ou generalização para outros serviços. Isso depende de revisão especializada de uma amostra aleatória, de leitura por mais de um examinador e de mais exemplos de lesões pequenas e maxilares.

\section*{Aprovação ética e consentimento}

\pendente{Informar o comitê de ética em pesquisa, o número do parecer (CAAE) e se houve consentimento dos pacientes ou dispensa formal, inclusive para a publicação das imagens.} As radiografias de pacientes foram processadas somente em computador local; nenhuma imagem de paciente foi enviada a serviços de nuvem. As figuras identificam pacientes apenas por números sequenciais.

\section*{Disponibilidade de dados e código}

O conjunto de Hanói está disponível publicamente [7]. As radiografias locais e as marcações não são públicas. \pendente{Informar a forma de acesso e as restrições.} O código de pré-processamento, treinamento, avaliação, figuras e revisão está no repositório do projeto \pendente{endereço do repositório}.

\section*{Contribuições dos autores}

\pendente{Descrever as contribuições segundo a taxonomia CRediT.}

\section*{Declaração de conflitos de interesse}

\pendente{Declarar conflitos de interesse ou a ausência deles.}

\section*{Uso de inteligência artificial}

Uma ferramenta de inteligência artificial generativa foi usada para implementar e executar as análises e para redigir este rascunho, a partir de dados, decisões e instruções dos autores. Todas as afirmações, números, referências e interpretações devem ser verificados e aprovados pelos autores, que permanecem responsáveis pelo conteúdo final.

\balance
\section*{Referências}

{\sffamily\fontsize{7.8pt}{10.2pt}\selectfont
\begin{enumerate}[leftmargin=5mm,itemsep=1mm]
\item Estrela, C., Bueno, M. R., Leles, C. R., Azevedo, B. \& Azevedo, J. R. Accuracy of cone beam computed tomography and panoramic and periapical radiography for detection of apical periodontitis. \textit{Journal of Endodontics} \textbf{34}, 273-279 (2008). \href{https://doi.org/10.1016/j.joen.2007.11.023}{doi:10.1016/j.joen.2007.11.023}
\item Endres, M. G. \textit{et al.} Development of a deep learning algorithm for periapical disease detection in dental radiographs. \textit{Diagnostics} \textbf{10}, 430 (2020). \href{https://doi.org/10.3390/diagnostics10060430}{doi:10.3390/diagnostics10060430}
\item Ekert, T. \textit{et al.} Deep learning for the radiographic detection of apical lesions. \textit{Journal of Endodontics} \textbf{45}, 917-922.e5 (2019).
\item Breiman, L. Random forests. \textit{Machine Learning} \textbf{45}, 5-32 (2001). \href{https://doi.org/10.1023/A:1010933404324}{doi:10.1023/A:1010933404324}
\item Berg, S. \textit{et al.} ilastik: interactive machine learning for (bio)image analysis. \textit{Nature Methods} \textbf{16}, 1226-1232 (2019). \href{https://doi.org/10.1038/s41592-019-0582-9}{doi:10.1038/s41592-019-0582-9}
\item Arganda-Carreras, I. \textit{et al.} Trainable Weka Segmentation: a machine learning tool for microscopy pixel classification. \textit{Bioinformatics} \textbf{33}, 2424-2426 (2017). \href{https://doi.org/10.1093/bioinformatics/btx180}{doi:10.1093/bioinformatics/btx180}
\item Do, H. V. \textit{et al.} A dataset of apical periodontitis lesions in panoramic radiographs for deep-learning-based classification and detection. \textit{Data in Brief} \textbf{54}, 110486 (2024). Dados: Panoramic radiographs with periapical lesions, Mendeley Data v3, \href{https://doi.org/10.17632/kx52tk2ddj.3}{doi:10.17632/kx52tk2ddj.3}
\item Hamamci, I. E. \textit{et al.} DENTEX: an abnormal tooth detection with dental enumeration and diagnosis benchmark for panoramic X-rays. arXiv:2305.19112 (2023).
\item Ronneberger, O., Fischer, P. \& Brox, T. U-Net: convolutional networks for biomedical image segmentation. In \textit{Medical Image Computing and Computer-Assisted Intervention -- MICCAI 2015}, LNCS 9351, 234-241 (2015). \href{https://doi.org/10.1007/978-3-319-24574-4_28}{doi:10.1007/978-3-319-24574-4\_28}
\item Ørstavik, D., Kerekes, K. \& Eriksen, H. M. The periapical index: a scoring system for radiographic assessment of apical periodontitis. \textit{Endodontics \& Dental Traumatology} \textbf{2}, 20-34 (1986). \href{https://doi.org/10.1111/j.1600-9657.1986.tb00119.x}{doi:10.1111/j.1600-9657.1986.tb00119.x}
\item Zuiderveld, K. Contrast limited adaptive histogram equalization. In \textit{Graphics Gems IV} (ed. Heckbert, P. S.) 474-485 (Academic Press, 1994).
\item He, K., Zhang, X., Ren, S. \& Sun, J. Deep residual learning for image recognition. In \textit{Proceedings of the IEEE Conference on Computer Vision and Pattern Recognition}, 770-778 (2016). \href{https://doi.org/10.1109/CVPR.2016.90}{doi:10.1109/CVPR.2016.90}
\item Loshchilov, I. \& Hutter, F. Decoupled weight decay regularization. In \textit{International Conference on Learning Representations} (2019).
\item Chakraborty, D. P. Maximum likelihood analysis of free-response receiver operating characteristic (FROC) data. \textit{Medical Physics} \textbf{16}, 561-568 (1989). \href{https://doi.org/10.1118/1.596358}{doi:10.1118/1.596358}
\item Wilson, E. B. Probable inference, the law of succession, and statistical inference. \textit{Journal of the American Statistical Association} \textbf{22}, 209-212 (1927). \href{https://doi.org/10.1080/01621459.1927.10502953}{doi:10.1080/01621459.1927.10502953}
\item Pedregosa, F. \textit{et al.} Scikit-learn: machine learning in Python. \textit{Journal of Machine Learning Research} \textbf{12}, 2825-2830 (2011).
\item Paszke, A. \textit{et al.} PyTorch: an imperative style, high-performance deep learning library. In \textit{Advances in Neural Information Processing Systems 32} (2019).
\item Salehi, S. S. M., Erdogmus, D. \& Gholipour, A. Tversky loss function for image segmentation using 3D fully convolutional deep networks. In \textit{Machine Learning in Medical Imaging}, LNCS 10541, 379-387 (2017). \href{https://doi.org/10.1007/978-3-319-67389-3_44}{doi:10.1007/978-3-319-67389-3\_44}
\item Ma, J. \textit{et al.} Segment anything in medical images. \textit{Nature Communications} \textbf{15}, 654 (2024). \href{https://doi.org/10.1038/s41467-024-44824-z}{doi:10.1038/s41467-024-44824-z}
\end{enumerate}}

\end{document}
